\SourcePage{331}{334}
At the opposite limiting extreme, the relativistic interaction is large relative to the electrostatic interaction (more precisely, relative to the part of the latter responsible for the energy dependence on $L$ and $S$). Orbital angular momentum and spin cannot then be considered separately, since neither is conserved. Each electron is characterized by its own total angular momentum $j$, and these momenta combine to give the total angular momentum $J$ of the atom. An atomic-level scheme of this kind is said to have $jj$ coupling. In fact, this coupling type does not occur in pure form; the levels of very heavy atoms exhibit various forms of coupling intermediate between the $LS$ and $jj$ types\footnote{For a fuller discussion of coupling types and their quantitative treatment, see E.~Condon and G.~Shortley, \emph{The Theory of Atomic Spectra}, Moscow: IL, 1949.}.
